\begin{abstract}
Video and EEG are widely used to study seizures in animals. However, manual
scoring remains a huge bottleneck. Making this automatic is a big challenge that
few have solved for animal models. This study tackles the problem using data from
twenty rats, analyzing about 24000 paired clips from 601 sessions. Experts graded
the seizures from a mild Stage\,2 to a severe Stage\,5. Ten video models and seven
EEG models were tested under strict rules to see how well they could spot seizures
and grade their exact severity. For basic detection, both tools did well, while
EEG was slightly better, scoring 0.974 compared to video's 0.968. Both tools also
worked well on unseen rats. The gap grew when grading the exact seizure levels. In
the five-level test, video scored 0.657, easily beating EEG's 0.560. The two
methods also made different types of mistakes. Video usually just confused stages
that were next to each other. EEG was more likely to miss seizures entirely or mix
up very different stages. The size of the video model did not matter much. A tiny
model did almost as well as a huge one. Instead, tracking motion over time was the
real key. An additional test on the outside RodEpil dataset confirmed these
results, with the model score reaching 0.979. Overall, we demonstrate that cameras
offer a great, pain-free way to track animal seizures and work very well alongside
typical EEG tools.
\end{abstract}
